\chapter{가지돌기 개수에 따른 뉴런}
\chaptermark{가지돌기 개수}
\label{sec:dendrites}

\section{회로 매개변수로서의 입력 개수}

\ref{sec:neurontypes}장에서는 뉴런을 출력이 무엇을 내보내는지에 따라, \ref{sec:eetypes}장에서는 어떤 소자로 동작하는지에 따라 분류했다. 엔지니어라면 곧바로 알아챌 세 번째 특징도 있다. 뉴런에 \emph{입력이 몇 개 있는가}이다. 입력은 가지돌기가 모은다. 가지돌기는 다른 세포의 축삭이 달라붙는 가지다(\ref{sec:dofa}장, \ref{sec:weightstore}절). 어떤 뉴런에는 가지돌기가 아예 없고, 어떤 뉴런에는 한두 개가 있으며, 또 어떤 뉴런에는 입력 10만 개를 모으는 나무가 있다. 전자 공학자에게 이것은 입력 부하 능력, 즉 팬인(fan-in)이며, 거의 언제나 과제에 맞춰져 있다.

가지돌기의 개수와 크기가 왜 그렇게 중요한지는 간단한 추정으로 알 수 있다. 막은 단위 면적당 커패시턴스가 $c_m\approx1$~\textmu F/cm$^2$인 커패시터이며, 가지돌기를 포함한 뉴런의 면적 $A$가 클수록 전압을 $\Delta V$만큼 올리는 데 더 많은 전하가 필요하다. 누설을 무시하면 입력 전류 $I$가 이 일을 하는 데 걸리는 시간은 다음과 같다.
\begin{empheq}[box=\keybox]{equation}
t \;\approx\; \frac{c_m\,A\,\Delta V}{I} .
\label{eq:dendriteload}
\end{empheq}
짧은 가지돌기가 몇 개 달린 작은 뉴런은 면적이 약 $2\cdot10^{-5}$~cm$^2$이고 커패시턴스는 약 $20$~pF이다. 전류가 몇 나노암페어인 큰 시냅스 하나가 이 뉴런을 10분의 1밀리초도 안 되어 $20\mV$ 올린다. 큰 나무를 가진 뉴런은 면적이 열다섯 배쯤 크고 커패시턴스는 약 $300$~pF이다. 전류가 약 $0.1$~nA인 보통 시냅스라면 이 뉴런을 충전하는 데 $60\ms$가 걸릴 것이다. 누설 시상수보다 훨씬 길므로, 결국 끝내 충전하지 못한다. 이런 뉴런에는 동시에 오는 입력 수십 개가 필요하다. 여기서 수치는 크기 수준일 뿐이지만 결론은 수치에 좌우되지 않는다. \emph{입력이 적으면 빠르고 정확하며, 입력이 많으면 느리고 합으로 동작한다}.

\begin{figure}[t]
\centering
\begin{tikzpicture}[>=Stealth,
  soma/.style={circle,draw,very thick,fill=blue!6,minimum size=0.55cm,inner sep=0pt},
  ax/.style={->,very thick,red!70!black},
  den/.style={very thick,blue!60!black}]
% 0 dendrites: sensor on a line
\begin{scope}[xshift=0cm]
\draw[den,decorate,decoration={snake,amplitude=1pt,segment length=4pt}] (-0.9,0) -- (0,0);
\node[font=\scriptsize] at (-0.9,0.3) {센서};
\draw[ax] (0,0) -- (1.0,0);
\draw[thick] (0.1,0) -- (0.1,0.45);
\node[soma,minimum size=0.4cm] at (0.1,0.65) {};
\node[font=\scriptsize,align=center] at (0.05,-0.75) {$0$\\센서 달린 선로};
\end{scope}
% 1 dendrite
\begin{scope}[xshift=3.3cm]
\draw[den] (-0.9,0) -- (-0.28,0);
\node[soma] at (0,0) {};
\draw[ax] (0.28,0) -- (0.9,0);
\node[font=\scriptsize,align=center] at (0,-0.75) {$1$\\버퍼, 리피터};
\end{scope}
% 2 dendrites: one per ear
\begin{scope}[xshift=6.2cm]
\draw[den] (-0.28,0) -- (-0.9,0); \draw[den] (0.28,0) -- (0.9,0);
\node[soma] at (0,0) {};
\draw[ax] (0,-0.28) -- (0,-0.6);
\node[font=\scriptsize] at (-0.9,0.25) {왼쪽}; \node[font=\scriptsize] at (0.9,0.25) {오른쪽};
\node[font=\scriptsize,align=center] at (0,-1.0) {$2$\\시간 AND};
\end{scope}
% ~4 short dendrites: mixer
\begin{scope}[xshift=9.0cm]
\foreach \a in {45,135,225,315}{\draw[den] (\a:0.28) -- (\a:0.7); \fill[blue!60!black] (\a:0.72) circle (0.06);}
\node[soma] at (0,0) {};
\draw[ax] (0.28,0) -- (1.0,0);
\node[font=\scriptsize,align=center] at (0,-1.0) {$\sim4$\\믹서};
\end{scope}
% many: summing tree
\begin{scope}[xshift=12.0cm]
\foreach \a in {60,90,120}{\draw[den] (0,0.28) -- ++(\a:0.55) coordinate (b); \foreach \d in {-20,20}{\draw[den] (b) -- ++(\a+\d:0.35);}}
\foreach \a in {-120,-60}{\draw[den] (0,-0.28) -- ++(\a:0.45);}
\node[soma] at (0,0) {};
\draw[ax] (0.28,0) -- (1.0,0);
\node[font=\scriptsize,align=center] at (0,-1.0) {$10^4$--$10^5$\\가산기};
\end{scope}
\end{tikzpicture}
\caption{입력 개수에 따른 다섯 부류. 파란색은 가지돌기(입력), 빨간색은 축삭(출력)이다. 입력이 적을수록 뉴런은 개별 신호를 더 빠르고 정확하게 전달하고, 많을수록 더 잘 더하고 분류한다. 세포 모양을 그린 것이 아니라 도식이다.}
\label{fig:dendriteclasses}
\end{figure}

\section{가지돌기 0개: 선로 끝의 센서}

촉각, 통증, 근육 신장을 감지하는 감각 뉴런(\ref{sec:sensors}, \ref{sec:pain}, \ref{sec:muscles}장)은 세포체에 가지돌기가 하나도 없다. 축삭은 세포체에서 가지 하나로 나와 곧 둘로 갈라진다. 한 가지는 피부나 근육으로 가고 그 끝이 곧 센서다. 다른 가지는 척수로 간다. 스파이크는 세포체를 거치지 않고 센서에서 곧장 척수로 달린다. 엔지니어에게 이것은 먼 끝에 센서가 달린 통신 선로이고, 세포체는 전원 장치처럼 옆에 매달려 있다. 세포체는 선로에 전력을 대지만 전송에는 참여하지 않는다. 이득은 속도와 신뢰성이다. 경로 어디에도 신호를 합산하고 전달할지 다시 결정해야 하는 곳이 없다.

빛과 소리의 센서는 더 극단적인 경우다. 이들은 입력을 모으는 뉴런이 아니라 변환기 자체이며, 입력은 시냅스가 아니라 수용체 단백질이다(그림~\ref{fig:transducer}).

\section{가지돌기 1개: 버퍼}

가지돌기 하나와 축삭 하나를 가진 뉴런은 입력 선로 하나를 받아 그대로 넘긴다. 후각 뉴런이 그렇다. 하나뿐인 가지돌기가 코의 표면으로 나가며 한 종류의 수용체만 지닌다~\cite{Mombaerts2004}. 망막에서 빛 센서와 눈의 출력 뉴런 사이에 있는 전달 세포, 귀의 센서를 뇌와 잇는 뉴런도 그렇다. 엔지니어라면 이들을 버퍼나 리피터라고 부를 것이다. 이들은 계산하지 않고 신호 하나를 배달하며, 때로는 \ref{sec:sensors}장에서처럼 센서의 매끄러운 전압을 스파이크로 바꾸는 식으로 신호의 형태를 바꾼다.

\section{가지돌기 2개: 시간 AND}

소리의 방향을 정하는 동시성 검출기(그림~\ref{fig:itd})는 포유류에서 흔히 서로 반대쪽을 향한 가지돌기를 정확히 두 개 가진다. 한쪽에는 왼쪽 귀의 입력이, 다른 쪽에는 오른쪽 귀의 입력이 온다~\cite{Grothe2010}. 왜 입력을 떼어 놓을까? 식~\eqref{eq:shunt}을 떠올려 보자. 막의 한 부위에서 흥분성 채널이 많이 열리면 그곳의 전압은 평형 전위에 가까워지고, 다음 입력은 점점 적게 보탠다. 그래서 한쪽 귀에서 온 입력이 한 가지돌기에 많이 모이면 그 가지돌기는 포화되어, 그것만으로는 뉴런을 문턱값까지 끌어올리지 못한다. 반면 각 가지돌기에서 하나씩 온 입력은 거의 선형으로 더해진다. 가지돌기 두 개는 뉴런을 더 엄격한 AND로 만든다. 신호가 양쪽에서 함께 왔을 때만 발화하라. 모델에서 이런 분리가 실제로 동시성 검출을 개선한다는 것이 밝혀졌다~\cite{AgmonSnir1998}.

\section{짧은 가지돌기 몇 개: 믹서와 중계기}

\emph{믹서.} \ref{sec:motorlearning}장의 운동 학습 회로에서는 약 700억 개의 작은 \nt{GLUT} 뉴런이 입력을 아주 긴 벡터로 확장한다. 이들 각각은 짧은 가지돌기가 약 네 개뿐이고, 가지돌기마다 입력 하나의 말단을 감싸는 “발톱”으로 끝난다. 따라서 믹서 하나는 많은 입력 가운데 약 네 개만 보고, 자기 네 입력에 대한 작은 AND 소자로 동작한다. 입력 선로 $M$개에서 서로 다른 네 개 묶음을 $\binom{M}{4}$가지 고를 수 있다. $M=100$이면 거의 400만 가지다. 왜 이렇게 많이 필요할까? 입력이 $n$개인 문턱 소자는 무작위 패턴을 최대
\begin{empheq}[box=\keybox]{equation}
P_{\max} \;\approx\; 2n
\label{eq:cover}
\end{empheq}
개까지 두 부류로 올바르게 나눌 수 있다~\cite{Cover1965}. 같은 회로의 출력 뉴런은 믹서에서 약 $10^5$개의 입력을 받으므로, \eqref{eq:cover}에 따라 상한은 서로 다른 대응 약 $2\cdot10^5$개다. 이것은 이상적인 소자의 한계다. 흥분성 시냅스처럼 모든 가중치가 양수이고 잡음에 대한 여유가 필요하면, 같은 뉴런에 대한 추정은 약 $4\cdot10^4$개의 대응이다~\cite{Brunel2004}. 입력이 적은 믹서는 바로 큰 가산기가 서로 다르고 거의 독립인 입력을 많이 갖게 하려고 필요하다~\cite{Marr1969,Albus1971}.

\emph{중계기.} 귀에서 방향 검출기로 가는 경로에는 짧은 가지돌기가 몇 개뿐인 뉴런이 있다. 이들은 거대한 시냅스를 몇 개만 받고, 어떤 뉴런은 사실상 하나만 받는다. \eqref{eq:dendriteload}에 따르면 이것이 딱 필요한 구성이다. 커패시턴스는 작고 전류는 크므로, 뉴런은 입력 스파이크마다 1밀리초도 안 되어 발화하며 시간 정밀도를 거의 잃지 않는다. 이런 중계기 가운데 하나는 \nt{GLUT}를 받아 \nt{GLYC}를 낸다. 수십 마이크로초 정밀도의 인버터로, 방향 검출기에 빠른 억제를 공급한다~\cite{BorstSoria2012}.

\section{입력 수천 개: 가산기와 분류기}

반대쪽 끝에는 입력이 1만 개쯤인 피질의 \nt{GLUT} 뉴런과, 입력이 10만 개인 운동 학습 회로의 출력 \nt{GABA} 뉴런이 있다. 이런 뉴런은 \eqref{eq:dendriteload}에 따라 느리고 개별 입력에 반응할 수 없다. 이들은 더한다. \ref{sec:glutcomputer}장의 적분 뉴런이며, 분류기를 만드는 바로 그 가산기다. 큰 나무는 새로운 것도 보탠다. 나무의 가지는 각자 문턱값을 가진 별도의 하위 회로로 동작하므로, 뉴런 하나가 작은 다층 네트워크처럼 행동한다~\cite{Beniaguev2021,Gidon2020}.

\section{출력이기도 한 가지돌기}

축삭이 아예 없고 가지돌기가 입력이자 출력인 뉴런도 있다. 이 뉴런은 신호를 받고 스스로 신경전달물질을 내보낸다. 가장 많이 연구된 예는 운동 방향 판별에 참여하는 망막의 \nt{GABA} 뉴런이다(\ref{sec:gaba}장, 그림~\ref{fig:direction}). 이런 뉴런에서는 가지 하나하나가 스스로 세포체에서 자기 끝 쪽으로 가는 움직임에 반응하고, 그 끝에서 억제를 낸다~\cite{Euler2002}. 뉴런 하나가 가지마다 하나씩, 독립된 방향 검출기 여러 개다. 엔지니어에게 이것은 소자 하나가 아니라 패키지 하나에 채널이 여럿 든 칩이다.

\section{요약}

\begin{table}[t]
\centering
\footnotesize
\setlength{\tabcolsep}{3pt}
\renewcommand{\arraystretch}{1.15}
\begin{tabular}{>{\raggedright}p{1.9cm}>{\raggedright}p{3.4cm}>{\raggedright}p{2.6cm}>{\raggedright}p{1.8cm}>{\raggedright\arraybackslash}p{3.8cm}}
\toprule
가지돌기 & 예 & 소자 & 입력 & 알려진 것 \\
\midrule
$0$ & 촉각, 통증, 신장 센서 & 끝에 센서가 달린 선로 & --- & 확립됨 \\
$1$ & 후각 뉴런, 망막의 전달 세포, 귀의 뉴런 & 버퍼, 리피터 & $1$--소수 & 확립됨 \\
$2$ & 소리 방향 검출기 & 시간 AND & 두 다발 & 구조는 확립됨; 입력 분리의 이득은 모델에서 입증 \\
$\sim4$ & 운동 학습 회로의 믹서 & 작은 AND 소자 & $\sim4$ & 확립됨; 입력 확장의 역할은 근거가 탄탄한 가설 \\
소수, 짧음 & 귀에서 오는 경로의 중계기 & 중계기, 인버터 & $1$--몇 개의 거대 시냅스 & 확립됨 \\
수천 & 피질의 \nt{GLUT} 뉴런, 운동 학습의 출력 뉴런 & 가산기, 분류기 & $10^4$--$10^5$ & 확립됨; 하위 회로로서의 가지는 개별 사례에서 측정 \\
출력 가지돌기 & 망막의 방향 \nt{GABA} 뉴런 & 다채널 칩 & 많음 & 측정됨 \\
\bottomrule
\end{tabular}
\caption{가지돌기 개수에 따른 뉴런.}
\label{tab:dendrites}
\end{table}

\begin{keyblock}
\begin{proposition}[입력 개수는 과제를 따른다]
\label{prop:dendrites}
센서, 중계기, 동시성 검출기처럼 개별 신호의 속도와 정밀도가 필요한 곳에서는 뉴런의 가지돌기가 적고 입력이 적으며 시냅스가 크다. 작은 커패시턴스~\eqref{eq:dendriteload}는 빨리 충전된다. 더하고 분류해야 하는 곳에서는 뉴런이 입력 수천 개를 가진 큰 나무를 갖는다. 부류의 구조는 측정되었다. 계산적 설명은 근거가 탄탄하지만 많은 부류에서는 아직 가설이다.
\end{proposition}
\end{keyblock}

표~\ref{tab:dendrites}는 앞의 두 분류, 즉 신경전달물질에 따른 분류(표~\ref{tab:types})와 소자에 따른 분류(표~\ref{tab:eetypes})를 보완한다. 세 분류 모두 같은 소자를 다른 쪽에서 본다. 무엇을 내보내는지, 무엇을 계산하는지, 얼마나 많이 듣는지다.

\section{연습문제}

\begin{exercise}[\lvlA]
\label{ex:19:1}
\emph{큰 뉴런에는 입력이 몇 개 필요한가.} $C=300$~pF, $\tau_m=20\ms$인 뉴런이 각각 $0.1$~nA의 전류원인 시냅스를 받고, 문턱값은 휴지 상태보다 $20\mV$ 높다. (a)~문턱값에 아예 도달하려면, $10\ms$ 안에, $5\ms$ 안에 도달하려면 시냅스 몇 개가 동시에 작동해야 하는가? (b)~$C=20$~pF인 뉴런은 $4$~nA 시냅스 하나로 얼마 만에 문턱값에 도달하는가?
\end{exercise}

\begin{exercise}[\lvlA]
\label{ex:19:2}
\emph{왜 네 개인가.} 믹서는 $M=100$개 가운데 $k$개 선로에 대한 AND이고, 패턴에서 선로의 절반이 활성이며, 믹서는 $7\cdot10^{10}$개다. (a)~$k=2$, $4$, $40$일 때 패턴에 반응하는 믹서는 몇 개인가? (b)~믹서는 직접 선로 $M=100$개에 비해 입력 $10^5$개인 출력 뉴런의 대응 수~\eqref{eq:cover}를 몇 배 늘리는가?
\end{exercise}

\begin{exercise}[\lvlB]
\label{ex:19:3}
\emph{$2n$은 어디서 오는가.} 원점을 지나는 초평면은 $n$차원 공간에서 일반 위치에 있는 점 $P$개를 $C(P,n)=2\sum_{k=0}^{n-1}\binom{P-1}{k}$가지 방식으로 두 부류로 나눈다. (a)~$P=2n$이면 $2^P$개 분할 가운데 정확히 절반이 분리 가능함을 증명하라. (b)~$n=100$, $P=180$과 $220$에서 분리 가능한 비율은 얼마인가?
\end{exercise}

\begin{exercise}[\lvlB]
\label{ex:19:4}
\emph{엄격한 AND로서의 가지돌기 두 개.} 가지돌기는 누설 $g_L$인 구획이고, 입력 하나가 그 위에 전위 $E_E$인 컨덕턴스 $g_0=0.5\,g_L$을 연다. 세포체는 $\kappa(V_1+V_2)$를 보고, 문턱값은 $\theta=1.1\,\kappa E_E$이다. 왜 한 가지돌기에 입력이 아무리 많아도 뉴런이 발화하지 않는가? 각 가지돌기에 입력이 몇 개 필요한가?
\end{exercise}

\begin{exercise}[\lvlB]
\label{ex:19:5}
\emph{중계기 설계.} 스파이크 지터는 $\sigma_t=\sigma_V/\dot V$이고, $\dot V$는 문턱값에서의 전압 기울기다. 잡음 $\sigma_V=1\mV$에서 $\sigma_t\le20$~\textmu s가 필요하다. 누설은 무시하라. $C=300$~pF인 뉴런에는 $0.1$~nA의 동기 시냅스가 몇 개 필요한가? $C=20$~pF이고 $4$~nA 시냅스가 하나인 뉴런의 지터는?
\end{exercise}

\begin{exercise}[\lvlB]
\label{ex:19:6}
\emph{모델링: 긴 가지돌기의 충전.} 끝이 닫힌 길이 $L=\ell/\lambda$의 수동 케이블 $\tau_m\partial_tV=\lambda^2\partial_x^2V-V+i/g$를 시뮬레이션하라($40$개 구획, 오일러법). 먼 끝에 짧은 전류 펄스를 넣는다. $L=0.2$, $1$, $2$에서 가까운 끝은 언제, 같은 전하를 가지돌기 전체에 고르게 퍼뜨렸을 때 전압의 몇 분의 몇까지 오르는가?
\end{exercise}

\begin{exercise}[\lvlC]
\label{ex:19:7}
\emph{탐구: 최적의 입력 개수.} 커패시턴스는 $C=C_0+nc_s$이고, 전류 $I_s$인 입력 $m$개가 동시에 작동하며, 뉴런은 대응 $P\approx2n$개를 기억한다~\eqref{eq:cover}. $m$이 일정할 때와 비율 $m=\varphi n$이 일정할 때 응답 시간은 $P$에 어떻게 의존하는가? 비용 기준을 제안하고 중계기와 분류기의 최적 $n$을 구하라.
\end{exercise}
